\section{The Lieb--Robinson input}\label{app:LR}

The only input from the theory of Lieb--Robinson bounds is the one-step estimate of \cite{NSY18} for lattice
fermions with even interactions, which we state in the special case needed here: a finite set of fermionic
modes and a time-independent interaction. In \cite{NSY18} each site carries one mode; spin is included by
taking as the set of sites the set $\Slots\times\{\up,\dn\}$ of modes, and the support of an operator in $\CAR(X)$,
$X\subset\Slots$, is then $X\times\{\up,\dn\}$.

Let $\mathfrak M$ be a finite set, $\CAR(Y)$ for $Y\subset\mathfrak M$ the algebra generated by the annihilation operators
$a_y$, $y\in Y$, and $\CAR(Y)^+$ its even part. An even interaction is a map $\Phi$ from subsets of $\mathfrak M$ to
$\CAR(\mathfrak M)$ with $\Phi(Z)=\Phi(Z)^*\in\CAR(Z)^+$; for $X\subset\mathfrak M$ let $H_X:=\sum_{Z\subset X}\Phi(Z)$ and
$\tau^X_t(A):=e^{\iu tH_X}Ae^{-\iu tH_X}$, and let $S(X):=\{Z\subset\mathfrak M:Z\cap X\ne\emptyset,\ Z\cap(\mathfrak M\setminus X)\ne\emptyset\}$.

\begin{proposition}[disjoint supports {\cite[Proposition~2.1]{NSY18}}]\label{prop:nsy-disjoint}
Let $X,Y\subset\mathfrak M$ be disjoint. If $A\in\CAR(X)^+$ and $B\in\CAR(Y)$, then $[A,B]=0$. If $A\in\CAR(X)$ and
$B\in\CAR(Y)$ are both odd, then $\{A,B\}=0$.
\end{proposition}

\begin{lemma}[{\cite[Lemma~3.2]{NSY18}}]\label{lem:nsy-ode}
Let $A(t)$, $B(t)$ be norm-continuous families of bounded operators on a Hilbert space, with $A(t)$
self-adjoint. Then the solution of $f'(t)=\iu[A(t),f(t)]+B(t)$, $f(t_0)=f_0$, satisfies
$\norm{f(t)}\le\norm{f(t_0)}+\int_{\min\{t,t_0\}}^{\max\{t,t_0\}}\norm{B(u)}\,\dd u$.
\end{lemma}

\begin{lemma}[the one-step estimate {\cite[Lemma~3.3]{NSY18}}]\label{lem:nsy-onestep}
Let $\Phi$ be an even interaction on $\mathfrak M$, $X\subset\mathfrak M$, $A\in\CAR(X)$, $B\in\CAR(\mathfrak M)$ and $t\ge0$. Then
\[
\bigl\lVert[\tau^{\mathfrak M}_t(A),B]\bigr\rVert\le\bigl\lVert[\tau^X_t(A),B]\bigr\rVert
+2\norm A\sum_{Z\in S(X)}\int_0^t\bigl\lVert[\tau^{\mathfrak M}_u(\Phi(Z)),B]\bigr\rVert\,\dd u .
\]
\end{lemma}

In \cite{NSY18} these statements are proved for time-dependent interactions and for all times $s\le t$; negative
times are reduced to positive ones in the proof below. Iterating
Lemma~\ref{lem:nsy-onestep} gives \cite[Theorem~3.1]{NSY18}, a Lieb--Robinson bound in terms of a decay function
on $\mathfrak M\times\mathfrak M$. Theorem~\ref{thm:LR} is obtained by the same iteration, with the sums over chains of
sets bounded by powers of the chain matrix instead of by convolutions of that decay function.

\subsection{Proof of Theorem~\ref{thm:LR}}

The bounds of Theorem~\ref{thm:LR} are not stated in \cite{NSY18}. We derive them here from
\cite[Lemma~3.3]{NSY18}, which is Lemma~\ref{lem:nsy-onestep}, by iterating it as in the proof of
\cite[Theorem~3.1]{NSY18}, and we bound the resulting sums over chains of sets by powers of the chain matrix where
\cite{NSY18} use the convolution property of a decay function. For a statement $S$ we write $[S]:=1$ if $S$ holds
and $[S]:=0$ otherwise.

\emph{The one-step estimate in terms of slots.} Let $\mathfrak M:=\Slots\times\{\up,\dn\}$ be the set of modes, with
$a_{(s,\sigma)}:=\chat_{s\sigma}$, and for $X\subset\Slots$ let $\widetilde X:=X\times\{\up,\dn\}$. The algebra $\CAR(X)$ of
Definition~\ref{def:slotfock} is the algebra $\CAR(\widetilde X)$ of this appendix, and $\Alg(X)=\CAR(\widetilde X)^+$: for
$A\in\CAR(\widetilde X)$, conjugation by $(-1)^{\Nhat\up+\Nhat\dn}$ coincides with conjugation by the parity of the modes in
$\widetilde X$, because the number operator of the remaining modes commutes with every $a_y$, $y\in\widetilde X$. By
\eqref{eq:loc-JW}, $\CAR(\mathfrak M)$ is the algebra of all operators on $\FockS$. Given the even interaction $\Psi$ on
$\Slots$, an even interaction $\Phi$ on $\mathfrak M$ in the sense of this appendix is defined by $\Phi(\widetilde Z):=\Psi(Z)$
for $Z\subset\Slots$ and $\Phi(W):=0$ for the subsets $W\subset\mathfrak M$ that are not of the form $\widetilde Z$. For
$X\subset\Slots$, $H_{\widetilde X}=\sum_{Z\subset X}\Psi(Z)=:H_X$; in particular $H_{\mathfrak M}=H_\Psi$. The sets of
$S(\widetilde X)$ on which $\Phi$ does not vanish are the $\widetilde Z$ with $Z$ in
\[
S_\Lambda(X):=\{Z\subset\Slots:Z\cap X\ne\emptyset,\ Z\not\subset X\}.
\]
Write $\tau^X_t$ for the dynamics generated by $H_X$, which is the dynamics $\tau^{\widetilde X}_t$ above; in particular
$\tau^{\Slots}_t=\tau^\Psi_t$. Then
Lemma~\ref{lem:nsy-onestep} says: for $X\subset\Slots$, $A\in\CAR(X)$, every operator $B$ on $\FockS$ and every $t\ge0$,
\begin{equation}\label{eq:LR-onestep}
\norm{[\tau^\Psi_t(A),B]}\le\norm{[\tau^X_t(A),B]}+2\norm A\sum_{Z\in S_\Lambda(X)}\int_0^t\norm{[\tau^\Psi_u(\Psi(Z)),B]}\,\dd u .
\end{equation}

\emph{Reduction to $t\ge0$.} For $t<0$, $\tau^\Psi_t=\tau^{-\Psi}_{\abs t}$, and $-\Psi$ is an even interaction with
$M^{-\Psi}=M^\Psi$ and $C(-\Psi,\lambda)=C(\Psi,\lambda)$. It therefore suffices to prove the theorem for $t\ge0$, which we
assume from now on.

\emph{The first term.} Let $W\subset\Slots$ and $u\ge0$. Since $H_W$ is an even element of $\CAR(W)$, conjugation by
$e^{\iu uH_W}$ maps $\CAR(W)$ onto itself and preserves parity, so $\tau^W_u$ maps $\Alg(W)$ onto itself. Hence, by
Proposition~\ref{prop:nsy-disjoint}, if either $A'\in\Alg(W)$ and $B\in\CAR(Y)$, or $A'\in\CAR(W)$ and $B\in\Alg(Y)$, then
$[\tau^W_u(A'),B]=0$ whenever $W\cap Y=\emptyset$; and $\norm{[\tau^W_u(A'),B]}\le2\norm{A'}\norm B$ in all cases.
Thus
\begin{equation}\label{eq:LR-first}
\norm{[\tau^W_u(A'),B]}\le2\norm{A'}\norm B\,[W\cap Y\ne\emptyset].
\end{equation}
Under either hypothesis of the theorem, \eqref{eq:LR-first} applies to $W=X$ and $A'=A$; and since every $\Psi(Z)$ is
even, it applies to $W=Z$ and $A'=\Psi(Z)$ for every $B\in\CAR(Y)$.

\emph{Iteration.} For $n\ge1$, a \emph{chain of length $n$ from $X$} is a sequence $(Z_1,\dots,Z_n)$ of subsets of
$\Slots$ with $Z_1\in S_\Lambda(X)$ and $Z_i\in S_\Lambda(Z_{i-1})$ for $2\le i\le n$. Put $a_0:=[X\cap Y\ne\emptyset]$ and,
for $n\ge1$,
\[
a_n:=\sum_{(Z_1,\dots,Z_n)}[Z_n\cap Y\ne\emptyset]\prod_{i=1}^n\norm{\Psi(Z_i)},\qquad
b_n:=\sum_{(Z_1,\dots,Z_n)}\prod_{i=1}^n\norm{\Psi(Z_i)},
\]
the sums running over the chains of length $n$ from $X$. We claim that for every $n\ge0$
\begin{equation}\label{eq:LR-iterate}
\norm{[\tau^\Psi_t(A),B]}\le2\norm A\norm B\sum_{j=0}^n\frac{(2t)^j}{j!}\,a_j+R_n(t),
\end{equation}
where, with the sum over the chains $(Z_1,\dots,Z_{n+1})$ of length $n+1$ from $X$,
\[
R_n(t):=2^{n+1}\norm A\sum_{(Z_1,\dots,Z_{n+1})}\prod_{i=1}^n\norm{\Psi(Z_i)}
\int_0^t\!\!\int_0^{u_1}\!\!\cdots\int_0^{u_n}\norm{[\tau^\Psi_{u_{n+1}}(\Psi(Z_{n+1})),B]}\,\dd u_{n+1}\cdots\dd u_1 .
\]
For $n=0$, \eqref{eq:LR-iterate} is \eqref{eq:LR-onestep} with its first term bounded by \eqref{eq:LR-first}. If
\eqref{eq:LR-iterate} holds for some $n$, apply \eqref{eq:LR-onestep} to $A'=\Psi(Z_{n+1})\in\CAR(Z_{n+1})$ at time
$u_{n+1}$, with its first term bounded by \eqref{eq:LR-first}:
\begin{multline*}
\norm{[\tau^\Psi_{u_{n+1}}(\Psi(Z_{n+1})),B]}\le2\norm{\Psi(Z_{n+1})}\norm B\,[Z_{n+1}\cap Y\ne\emptyset]\\
+2\norm{\Psi(Z_{n+1})}\sum_{Z_{n+2}\in S_\Lambda(Z_{n+1})}\int_0^{u_{n+1}}\norm{[\tau^\Psi_{u_{n+2}}(\Psi(Z_{n+2})),B]}\,\dd u_{n+2}.
\end{multline*}
Inserted in $R_n(t)$, the first term contributes $2^{n+1}\norm A\cdot2\norm B\,a_{n+1}\,t^{n+1}/(n+1)!$, because the
simplex $\{t\ge u_1\ge\cdots\ge u_{n+1}\ge0\}$ has volume $t^{n+1}/(n+1)!$; this is the term $j=n+1$ of
\eqref{eq:LR-iterate}. The second term contributes exactly $R_{n+1}(t)$. This proves \eqref{eq:LR-iterate} for every $n$.
Bounding the integrand of $R_n(t)$ by $2\norm{\Psi(Z_{n+1})}\norm B$ gives
\[
R_n(t)\le2\norm A\norm B\,\frac{(2t)^{n+1}}{(n+1)!}\,b_{n+1}.
\]

\emph{Chains and powers of the chain matrix.} Let $M:=M^\Psi$ and $\mu:=\max_s\sum_{s'}M_{ss'}$. We claim that for every
$n\ge0$
\[
a_n\le\sum_{s\in X,\ s'\in Y}(M^n)_{ss'},\qquad b_n\le\sum_{s\in X,\ s'\in\Slots}(M^n)_{ss'}\le\abs X\mu^n .
\]
For $n=0$, $a_0\le\abs{X\cap Y}=\sum_{s\in X,\,s'\in Y}[s=s']$. Let $n\ge1$, and let $(Z_1,\dots,Z_n)$ be a chain from $X$
with $Z_n\cap Y\ne\emptyset$. Choose $s_1\in Z_1\cap X$, $s_i\in Z_{i-1}\cap Z_i$ for $2\le i\le n$ (these intersections are
nonempty because $Z_i\in S_\Lambda(Z_{i-1})$), and $s_{n+1}\in Z_n\cap Y$. Then $s_i,s_{i+1}\in Z_i$ for $1\le i\le n$, so the
term $\prod_i\norm{\Psi(Z_i)}$ of the chain is one of the nonnegative terms obtained by expanding
\[
\sum_{s_1\in X}\ \sum_{s_2,\dots,s_n\in\Slots}\ \sum_{s_{n+1}\in Y}\ \prod_{i=1}^n\Bigl(\sum_{Z_i\ni s_i,s_{i+1}}\norm{\Psi(Z_i)}\Bigr)
=\sum_{s\in X,\ s'\in Y}(M^n)_{ss'},
\]
and distinct chains give distinct terms. This proves the bound on $a_n$. The bound on $b_n$ follows in the same way,
with $s_{n+1}\in Z_n$ arbitrary, and $\sum_{s'}(M^n)_{ss'}\le\mu^n$.

\emph{The first bound.} By the last two steps, $R_n(t)\le2\norm A\norm B\abs X(2t\mu)^{n+1}/(n+1)!\to0$ as $n\to\infty$.
Letting $n\to\infty$ in \eqref{eq:LR-iterate},
\[
\norm{[\tau^\Psi_t(A),B]}\le2\norm A\norm B\sum_{j\ge0}\frac{(2t)^j}{j!}\,a_j\le2\norm A\norm B\sum_{s\in X,\ s'\in Y}\bigl[e^{2tM}\bigr]_{ss'}.
\]
If $X\cap Y=\emptyset$, then $a_0=0$, and the term $j=0$, which contributes the identity matrix, can be omitted; this
replaces $e^{2tM}$ by $e^{2tM}-\one$.

\emph{The weighted bound.} Let $\lambda>0$, $C:=C(\Psi,\lambda)$ and $M^{\rm w}_{ss'}:=M_{ss'}e^{d(s,s')/\lambda}$, so that
every row of $M^{\rm w}$ sums to at most $C$. For slots $s_1,\dots,s_{n+1}$ the triangle inequality
$d(s_1,s_{n+1})\le\sum_id(s_i,s_{i+1})$ gives $\prod_{i=1}^nM_{s_is_{i+1}}\le e^{-d(s_1,s_{n+1})/\lambda}\prod_{i=1}^nM^{\rm w}_{s_is_{i+1}}$;
summing over $s_2,\dots,s_n$, $(M^n)_{ss'}\le e^{-d(s,s')/\lambda}\bigl((M^{\rm w})^n\bigr)_{ss'}$. Hence, for every slot $s$,
every $W\subset\Slots$ and every $n\ge1$,
$\sum_{s'\in W}(M^n)_{ss'}\le e^{-d(s,W)/\lambda}\sum_{s'\in\Slots}((M^{\rm w})^n)_{ss'}\le e^{-d(s,W)/\lambda}C^n$, while
$\sum_{s'\in W}[s=s']=[s\in W]\le e^{-d(s,W)/\lambda}$. Summing the exponential series, for every $t\in\R$,
\begin{equation}\label{eq:LR-rowsum}
\sum_{s'\in W}\bigl[e^{2\abs tM}\bigr]_{ss'}\le e^{2\abs tC}e^{-d(s,W)/\lambda},\qquad
\sum_{s'\in W}\bigl[e^{2\abs tM}-\one\bigr]_{ss'}\le\bigl(e^{2\abs tC}-1\bigr)e^{-d(s,W)/\lambda}.
\end{equation}
Summing \eqref{eq:LR-rowsum} with $W=Y$ over $s\in X$ and inserting it in the first bound gives the weighted bound,
with $e^{2\abs tC}-1$ in place of $e^{2\abs tC}$ when $X\cap Y=\emptyset$.

\emph{Deleting terms.} If $\Psi'$ is obtained from $\Psi$ by setting some of its terms to zero, then
$0\le M^{\Psi'}\le M^\Psi$ entrywise, hence $(M^{\Psi'})^n\le(M^\Psi)^n$ entrywise for every $n$, and
$C(\Psi',\lambda)\le C(\Psi,\lambda)$. The right-hand sides of both bounds, and of \eqref{eq:LR-rowsum}, are nondecreasing
functions of the entries of the matrix, so the bounds for the dynamics of $\Psi'$ hold with the matrix $M^\Psi$.
\qed
